\documentclass[12pt,letterpaper]{article}

% Compile from this directory: latexmk -pdf -outdir=build writeup.tex
\usepackage[T1]{fontenc}
\usepackage{lmodern}
\usepackage[margin=1in]{geometry}
\usepackage{microtype}
\usepackage{amsmath,amssymb,amsthm,mathtools}
\usepackage{enumitem}
\usepackage{xltabular}
\usepackage{graphicx}
\usepackage{tikz-cd}
\usepackage{aliascnt}
\usepackage{xcolor}
\definecolor{accent}{HTML}{176B70}
\usepackage[colorlinks=true,linkcolor=accent,citecolor=accent,urlcolor=accent]{hyperref}
\usepackage[nameinlink,noabbrev]{cleveref}
\usepackage{mathpazo}

% Gaps scale with font size, independently of font weight.
% Number-to-title and punctuation spacing both use half the font size.
\makeatletter
% Keep the font-size calculation unexpanded while LaTeX prepares section labels.
\protected\def\headingnumbersep{\dimexpr\f@size pt/2\relax}
\protected\def\headingpunctsep{\dimexpr\f@size pt/2\relax}
\renewcommand{\@seccntformat}[1]{%
  \csname the#1\endcsname\hspace{\headingnumbersep}%
}
\makeatother

\numberwithin{equation}{subsection}
% One sequence per subsection: section.subsection.index.
% The period belongs after the environment title, before the optional name.
\newtheoremstyle{numberfirstplain}
  {\topsep}{\topsep}{\itshape}{}{\bfseries}{}{\headingpunctsep}
  {\thmnumber{#2}\thmname{\hspace{\headingnumbersep}#1.}\thmnote{\hspace{\headingpunctsep}{\normalfont(#3)}}}
\newtheoremstyle{numberfirstupright}
  {\topsep}{\topsep}{\normalfont}{}{\bfseries}{}{\headingpunctsep}
  {\thmnumber{#2}\thmname{\hspace{\headingnumbersep}#1.}\thmnote{\hspace{\headingpunctsep}{\normalfont(#3)}}}
\theoremstyle{numberfirstplain}
\newtheorem{theorem}{Theorem}[subsection]
% Alias counters preserve each environment's name in \cref references.
\newcommand{\sharedtheorem}[2]{%
  \newaliascnt{#1}{theorem}%
  \newtheorem{#1}[#1]{#2}%
  \aliascntresetthe{#1}%
}
\sharedtheorem{proposition}{Proposition}
\sharedtheorem{lemma}{Lemma}
\sharedtheorem{corollary}{Corollary}
\theoremstyle{numberfirstupright}
\sharedtheorem{definition}{Definition}
\sharedtheorem{notation}{Notation}
\sharedtheorem{example}{Example}
\sharedtheorem{question}{Question}
\sharedtheorem{remark}{Remark}
% Usage: \begin{custom}[Optional Info]{Your heading} ... \end{custom}
% Headings may contain linked references, e.g. \Cref{thm:main}.
\newcommand{\customheading}{}
\sharedtheorem{custom}{\customheading}
\let\customstart\custom
\let\customfinish\endcustom
\RenewDocumentEnvironment{custom}{O{} m}
  {\renewcommand{\customheading}{#2}\customstart[#1]}
  {\customfinish}
\crefname{custom}{item}{items}
\Crefname{custom}{Item}{Items}
\crefname{question}{question}{questions}
\Crefname{question}{Question}{Questions}
\crefname{notation}{notation}{notations}
\Crefname{notation}{Notation}{Notations}
% Section references use section symbols; nameinlink links symbol and number.
\crefname{section}{\S}{\S\S}
\Crefname{section}{\S}{\S\S}
\crefname{subsection}{\S}{\S\S}
\Crefname{subsection}{\S}{\S\S}
\crefname{subsubsection}{\S}{\S\S}
\Crefname{subsubsection}{\S}{\S\S}

% Centered, page-breakable tables, using their natural width up to the width limit.
% notationtable rows: symbol & type & description \\
% Declarations have aligned colons and no surrounding brackets.
% typetable rows: type & description \\
% Mathematical columns are automatically in math mode and use their natural
% widths without wrapping; only the rightmost description column wraps.
\newcommand{\notationtablespace}{0.4\baselineskip}
\newcommand{\notationtablewidth}{\linewidth}
\newsavebox{\notationtablemeasure}
% Emit a pending theorem/list heading before starting a vertical table.
\makeatletter
\newcommand{\notationtableflushheading}{\if@inlabel\leavevmode\fi}
\makeatother
\NewDocumentEnvironment{notationtable}{+b}{%
  \notationtableflushheading
  \par\nobreak\addvspace{\notationtablespace}
  \begingroup
  \setlength{\LTleft}{0pt plus 1fill}%
  \setlength{\LTright}{0pt plus 1fill}%
  \setlength{\LTpre}{0pt}%
  \setlength{\LTpost}{0pt}%
  \renewcommand{\arraystretch}{1.2}%
  \sbox{\notationtablemeasure}{%
    \begin{tabular}{@{}>{$}r<{$}@{\hspace{0.4em}:\hspace{0.4em}}>{$}l<{$}@{\hspace{1.2em}}l@{}}
      #1
    \end{tabular}%
  }%
  \ifdim\wd\notationtablemeasure>\notationtablewidth\relax
    \begin{xltabular}{\notationtablewidth}{@{}>{$}r<{$}@{\hspace{0.4em}:\hspace{0.4em}}>{$}l<{$}@{\hspace{1.2em}}>{\raggedright\arraybackslash}X@{}}
      #1
    \end{xltabular}%
  \else
    \begin{longtable}{@{}>{$}r<{$}@{\hspace{0.4em}:\hspace{0.4em}}>{$}l<{$}@{\hspace{1.2em}}l@{}}
      #1
    \end{longtable}%
  \fi\par
  \endgroup
  \addvspace{\notationtablespace}
}{}
\NewDocumentEnvironment{typetable}{+b}{%
  \notationtableflushheading
  \par\nobreak\addvspace{\notationtablespace}
  \begingroup
  \setlength{\LTleft}{0pt plus 1fill}%
  \setlength{\LTright}{0pt plus 1fill}%
  \setlength{\LTpre}{0pt}%
  \setlength{\LTpost}{0pt}%
  \renewcommand{\arraystretch}{1.2}%
  \sbox{\notationtablemeasure}{%
    \begin{tabular}{@{}>{$}c<{$}@{\hspace{1.2em}}l@{}}
      #1
    \end{tabular}%
  }%
  \ifdim\wd\notationtablemeasure>\notationtablewidth\relax
    \begin{xltabular}{\notationtablewidth}{@{}>{$}c<{$}@{\hspace{1.2em}}>{\raggedright\arraybackslash}X@{}}
      #1
    \end{xltabular}%
  \else
    \begin{longtable}{@{}>{$}c<{$}@{\hspace{1.2em}}l@{}}
      #1
    \end{longtable}%
  \fi\par
  \endgroup
  \addvspace{\notationtablespace}
}{}

% \paragraph{} Text produces "2.1.3. Text" and shares the theorem counter.
% A nonempty argument adds a bold run-in title. \paragraph* stays unnumbered.
% Ordinary body paragraphs are also numbered automatically (see hooks below).
\newaliascnt{numberedparagraph}{theorem}
\crefname{numberedparagraph}{paragraph}{paragraphs}
\Crefname{numberedparagraph}{Paragraph}{Paragraphs}
% Ignore source whitespace even when one or more labels precede the body.
% A plain \ignorespaces stops at \label, leaving its following newline visible.
\makeatletter
\newcommand{\numberedparagraphignorespaces}{%
  \@ifnextchar\label{\numberedparagraphlabel}{\ignorespaces}%
}
\def\numberedparagraphlabel\label{\numberedparagraphlabelargs}
\NewDocumentCommand{\numberedparagraphlabelargs}{o m}{%
  \IfNoValueTF{#1}{\label{#2}}{\label[#1]{#2}}%
  \numberedparagraphignorespaces
}
\makeatother
\let\originalparagraph\paragraph
\RenewDocumentCommand{\paragraph}{s m}{%
  \IfBooleanTF{#1}{\SkipAutomaticParagraphNumber\originalparagraph*{#2}}{%
    \par\addvspace{\topsep}%
    \refstepcounter{numberedparagraph}%
    \SkipAutomaticParagraphNumber
    \noindent\textbf{\thenumberedparagraph.\hspace{\headingpunctsep}}%
    \if\relax\detokenize{#2}\relax\else\textbf{#2.}\hspace{\headingpunctsep}\fi
    \numberedparagraphignorespaces
  }%
}

% Number top-level prose only after a source paragraph break (a blank line
% or an explicit \par). Layout paragraphs inside environments do not arm
% numbering, so prose immediately following a table continues the same block.
\makeatletter
\newcount\autopar@suspended
\newif\ifautopar@skipnext
\newif\ifautopar@pending
\newcommand{\SkipAutomaticParagraphNumber}{\global\autopar@skipnexttrue}
\def\autopar@document{document}
\newcommand{\autopar@sourcebreak}{%
  \ifnum\autopar@suspended=0
    \ifx\@currenvir\autopar@document
      \ifinner\else
        \global\autopar@pendingtrue
      \fi
    \fi
  \fi
}
% Preserve LaTeX's paragraph machinery, including its environment-restoration
% path through \@@par. A blank line in vertical mode must also be detected.
\AddToHook{begindocument/end}{%
  \let\autopar@originalpar\par
  \protected\def\par{%
    \autopar@sourcebreak\autopar@originalpar\autopar@blockspace}%
  \let\@@par\par
}
% Add theorem-style separation at the source break, before TeX inserts the
% next paragraph's glue. Adjacent theorem skips then merge via \addvspace.
\newcommand{\autopar@blockspace}{%
  \ifnum\autopar@suspended=0
    \ifnum\value{section}>0
      \ifx\@currenvir\autopar@document
        \ifinner\else
          \ifautopar@skipnext\else
            \ifautopar@pending
              \if@nobreak\else
                \addvspace{\topsep}%
              \fi
            \fi
          \fi
        \fi
      \fi
    \fi
  \fi
}
\AddToHook{para/begin}[candidacy/numbered-prose]{%
  \ifnum\autopar@suspended=0
    \ifnum\value{section}>0
      \ifinner\else
        \ifx\@currenvir\autopar@document
          \ifautopar@skipnext
            \global\autopar@skipnextfalse
          \else
            \ifautopar@pending
              \OmitIndent
              \refstepcounter{numberedparagraph}%
              \textbf{\thenumberedparagraph.\hspace{\headingpunctsep}}%
            \fi
          \fi
        \fi
        % An environment's body also consumes a preceding source break.
        \global\autopar@pendingfalse
      \fi
    \fi
  \fi
}
\newcommand{\autopar@excludecommand}[1]{%
  \AddToHook{cmd/#1/before}{%
    \global\autopar@pendingfalse
    \advance\autopar@suspended by 1\relax}%
  \AddToHook{cmd/#1/after}{\advance\autopar@suspended by -1\relax}%
}
% Hook the fully parsed heading commands, not argument-scanning \section.
\autopar@excludecommand{@sect}
\autopar@excludecommand{@ssect}
\autopar@excludecommand{maketitle}
\autopar@excludecommand{tableofcontents}
\autopar@excludecommand{@footnotetext}
\autopar@excludecommand{@mpfootnotetext}
\autopar@excludecommand{newpage}
\autopar@excludecommand{clearpage}
\makeatother
% Use at the start of ordinary prose when a paragraph needs a reference label.
\NewDocumentCommand{\parlabel}{m}{\leavevmode\label{#1}\ignorespaces}
\NewDocumentEnvironment{unnumberedparagraphs}{}{\par}{\par}
\renewcommand{\refname}{References}

% Notation: keep consistent with presentation/presentation.tex.
\newcommand{\C}{\mathbb{C}}
\newcommand{\Z}{\mathbb{Z}}
\newcommand{\PP}{\mathbb{P}}
\DeclareMathOperator{\GL}{GL}
\DeclareMathOperator{\End}{End}
\DeclareMathOperator{\Hom}{Hom}
\DeclareMathOperator{\rank}{rank}
\DeclareMathOperator{\rig}{rig}

\title{Middle Convolution of $\ell$-adic Perverse Sheaves\\[0.4em]\small Submitted for the Ph.D. Candidacy Examination at Western University}
\author{Yunhai Xiang}
\date{\today}

\begin{document}

\maketitle
\tableofcontents
\newpage

\section{Introduction}\label{sec:introduction}


\subsection{Notations}\label{subsec:notation-style}


Motivated by the growing interest in the formalization of mathematics within the mathematical community, we adopt a style of notation inspired by type theory throughout this write-up, without committing to type-theoretic foundations. The symbols below denote types and/or categories depending on your philosophical inclination, where \(X:\mathsf{T}\) denotes that \(X\) is a term of type \(\mathsf{T}\), or an object of the category \(\mathsf{T}\).
\begin{typetable}
  \mathsf{Set} & Sets \\
  \mathsf{Set}^{\mathrm{fin}} & Finite sets \\
  \mathsf{Top} & Topological spaces \\
  \mathsf{Grp} & Groups \\
  \mathsf{Grp}^{\mathrm{top}} & Topological groups \\
  \mathsf{Grp}^{\mathrm{pro}} & Profinite groups \\
  \mathsf{Grp}^{\mathrm{ab}} & Abelian groups \\
  \mathsf{Ring} & Rings (associative, unital, not necessarily commutative) \\
  \mathsf{Ring}^{\mathrm{top}} & Topological rings \\
  \mathsf{Field} & Fields \\
  \mathsf{Field}^{\mathrm{top}} & Topological fields \\
  \mathsf{Mod}_R & Left modules over \(R : \mathsf{Ring}\) \\
  \mathsf{Vect}_k & Vector spaces over \(k : \mathsf{Field}\) \\
  \mathsf{Vect}^{\mathrm{fin}}_k & Finite-dimensional vector spaces over \(k : \mathsf{Field}\) \\
  \mathsf{Rep}_\Lambda(G) & Continuous \(\Lambda\)-linear representations of \(G : \mathsf{Grp}^{\mathrm{top}}\), where \(\Lambda : \mathsf{Ring}^{\mathrm{top}}\) \\
  \mathsf{Rep}^{\mathrm{fin}}_\Lambda(G) & Continuous \(\Lambda\)-linear representations of \(G : \mathsf{Grp}^{\mathrm{top}}\), where \(\Lambda : \mathsf{Ring}^{\mathrm{top}}\), whose underlying topological \(\Lambda\)-modules are finitely generated \\
  \mathsf{Sch} & Schemes \\
  \mathsf{Sch}_S & Schemes over base scheme $S:\mathsf{Sch}$\\
  \mathsf{Var}_k & Varieties over \(k : \mathsf{Field}\), i.e. integral separated schemes of finite type over \(k\) \\
  \mathsf{Curv}_k & Curves over \(k : \mathsf{Field}\), i.e. one-dimensional varieties over \(k\) \\
  \mathsf{Curv}^{\mathrm{sm}}_k & Smooth curves over \(k : \mathsf{Field}\) \\
  \mathsf{Man}_{\C} & Complex manifolds \\
  \mathsf{Riem} & Riemann surfaces, i.e. connected one-dimensional complex manifolds \\
  \mathsf{Site} & Sites, i.e. categories equipped with Grothendieck topologies \\
  \mathsf{Sh}(X,\mathsf A) & Sheaves on \(X : \mathsf{Site}\) with values in a category \(\mathsf A\) admitting small limits\\
  \mathsf{Mod}(X,\mathcal R) & Sheaves of left \(\mathcal R\)-modules on \(X : \mathsf{Site}\), where \(\mathcal R : \mathsf{Sh}(X,\mathsf{Ring})\) \\
  \mathsf D^{\mathrm b}(\mathsf A) & Bounded derived objects of an abelian category \(\mathsf A\) \\
  \mathsf D^{\mathrm b}(X,k) & Bounded derived objects of \(\mathsf{Sh}(X,\mathsf{Vect}_k)\), where \(X : \mathsf{Site}\) and \(k : \mathsf{Field}\) \\
  X\rightarrow Y & Morphisms from $X : \mathsf{T}$ to $Y : \mathsf{T}$
\end{typetable}

\begin{custom}{Global Notations}Let \(X,Y : \mathsf{Site}\), \(k : \mathsf{Field}\), \(n : \mathbb{Z}\), \(V : \mathsf{Vect}_k\), and \(f : X\to Y\). Also let \(T : \mathsf{Top}\) with base point \(t : T\). The following notations will be used throughout the write-up.
\begin{notationtable}
  (-)^{\mathrm{an}} & \mathsf{Curv}^{\mathrm{sm}}_{\C}\to\mathsf{Riem} & The complex analytification functor \\
  {[n]} & \mathsf D^{\mathrm b}(X,k)\to\mathsf D^{\mathrm b}(X,k) & The cohomological shift functor by \(n\)\\
  \Gamma(X,-) & \mathsf D^{\mathrm b}(X,k)\to\mathsf D^{\mathrm b}(\mathsf{Vect}_k) & The derived global sections functor\\
  \mathcal{H}^n & \mathsf D^{\mathrm b}(X,k)\to\mathsf{Sh}(X,\mathsf{Vect}_k) & The \(n\)-th cohomology sheaf functor \\
  \mathrm{H}^n(X,-) & \mathsf D^{\mathrm b}(X,k)\to\mathsf{Vect}_k & The \(n\)-th hypercohomology functor\\
  \underline{V} & \mathsf{Sh}(X,\mathsf{Vect}_k) & The constant sheaf associated with \(V\) \\
  f_* & \mathsf D^{\mathrm b}(X,k)\to\mathsf D^{\mathrm b}(Y,k) & The derived pushforward along \(f\) \\
  f^* & \mathsf D^{\mathrm b}(Y,k)\to\mathsf D^{\mathrm b}(X,k) & The derived pullback along \(f\)\\
  \otimes_k & \mathsf D^{\mathrm b}(X,k)\times\mathsf D^{\mathrm b}(X,k)\to\mathsf D^{\mathrm b}(X,k) & The derived tensor product \\
  \mathcal{H}om & \mathsf D^{\mathrm b}(X,k)^{\mathrm{op}}\times\mathsf D^{\mathrm b}(X,k)\to\mathsf D^{\mathrm b}(X,k) & The derived sheaf Hom\\
  \mathcal{E}nd & \mathsf D^{\mathrm b}(X,k)\to\mathsf D^{\mathrm b}(X,k) & The derived sheaf End \\
  \pi_1^{\mathrm{top}}(T,t) & \mathsf{Grp}^{\mathrm{top}} & The topological fundamental group, equipped with the discrete topology \\
\end{notationtable}
\noindent where the derived functors are assumed to preserve boundedness.
\end{custom}
\subsection{Motivations}\label{subsec:motivations}

\textit{Monodromy groups} are interesting mathematical objects that arise naturally in both geometric and arithmetic contexts. 
Roughly speaking, they describe how locally defined data change when transported along a loop. 
Controlling and understanding monodromy groups has been a central theme in both of these worlds, with significant applications.
We will discuss both the geometric and arithmetic sides, but this write-up will primarily focus on the arithmetic side, with the geometric side serving as motivations and intuitions. 

Let $w:\mathbb C$ be nonzero, and consider the equation
$$z^2=w$$
Locally near each $w$, there are two solutions which we label as $z_{+}=\sqrt{w}$ and $z_{-}=-\sqrt{w}$. Suppose we move $w$ around the unit circle with 
$$w(t)=e^{2\pi i t}\quad t:\mathbb R$$
Then we have $z_{+}(t)=e^{\pi i t}$ and $z_{-}(t)=-e^{\pi i t}$. When $t$ goes from $0$ to $1$, the two solutions do not return to their original self, but swap places. 
This induces a permutation $\mathrm{S}_2$ on the set of solutions $\{z_{+},z_{-}\}$. In other words, there is a permutation representation of the fundamental group $\pi^{\mathrm{top}}_1(\mathrm S^1,1)\cong\mathbb Z$
$$\pi^{\mathrm{top}}_1(\mathrm S^1,1)\longrightarrow\mathrm{Aut}_{\mathsf{Set}}(\{z_+,z_-\})\quad n\mapsto (z_{+}\ z_{-})^{(-1)^{n+1}}$$
We call this permutation representation the \textit{permutation monodromy representation}, and the image of this representation $\mathrm{S}_2$ the \textit{permutation monodromy group}, of the covering map 
$$p:\mathbb C^\times\rightarrow\mathbb C^\times\quad z\mapsto z^2$$

We can upgrade permutation monodromy representations into linear representations. 

\begin{custom}[]{A Question in Arithmetic Geometry}
Here is a 
\end{custom}




\subsection{Structure of the Write-up}\label{subsec:structure}

In \cref{sec:local-systems}

In \cref{sec:perverse-sheaves}

In \cref{sec:middle-convolution}

In \cref{sec:hyperelliptic-monodromy}

In \cref{sec:formalization}

\subsection{Acknowledgements}\label{subsec:acknowledgements}

The author thanks his advisor, Prof. Chris Hall, for his guidance and advice.

\begin{custom}[]{Typesetting}
The typesetting style of this write-up is inspired by \cite{Katz1988GKM}.
\end{custom}


\begin{custom}[]{Declaration of AI Use}
The use of generative AI in the preparation of this write-up was limited to the following purposes: 
polishing wording, phrasing, spelling, grammar, and formatting; 
offering writing advices; 
answering mathematical questions from the author to help him learn the subject;
searching and citing appropriate references; 
translating reading materials in French into English; 
retypesetting old references with archaic typesettings for ease of reading;
providing emotional support for the author;
providing fitness advice for the author to stay healthy and energetic while working on the write-up. 
\end{custom}



\section{Local Systems and Lisse Sheaves}\label{sec:local-systems}


\subsection{Local Systems in Topology}

\begin{notation}Define the following collection of notations
\begin{notationtable}
  X & \mathsf{Curv}^{\mathrm{sm}}_{\C} & A smooth projective connected curve over \(\C\) \\
  S & \mathsf{Set}^{\mathrm{fin}} & A finite nonempty subset of closed points of \(X\) \\
  U & \mathsf{Curv}^{\mathrm{sm}}_{\C} & The open subvariety \(X\setminus S\)\\
  j & U\rightarrow X & The inclusion map \\
  m & \mathbb{N} & The cardinality of \(S\)
\end{notationtable}
\end{notation}

\subsection{Local Systems in \texorpdfstring{$\ell$}{l}-adic Arithmetic Geometry}


\section{Perverse Sheaves and the Six Functor Formalism}\label{sec:perverse-sheaves}
\section{Middle Convolution and Monodromy Groups}\label{sec:middle-convolution}
\section{Mod-\texorpdfstring{$\ell$}{l} Monodromy of Hyperelliptic Curves}\label{sec:hyperelliptic-monodromy}

\section{Formalization in Lean}\label{sec:formalization}

% Bibliography entries live in references.bib.
\bibliographystyle{alpha}
\bibliography{references}

\end{document}
